\section*{Chapter \ref{ch:s1:cross-asset-lead-lag} --- Cross-Asset Lead--Lag}

\begin{solution}{exo:s1:cross-asset-lead-lag:1}
$\Phi^{-1}(1 - 0.05/2\,000) = 4.06$ for 1\,000 tests; $\Phi^{-1}(1 - 0.05/200) = 3.48$ for 100.
\end{solution}

\begin{solution}{exo:s1:cross-asset-lead-lag:2}
At one lag the covariance is $0.08\,\sigma_x^2/5$, so the correlation is $0.08 \times (25/8)/5 = 0.05$ and $t = 0.05\sqrt{1\,260} = 1.77$. With the
five-day window the covariance is $0.08\,\sigma_x^2$ and the window's volatility $\sigma_x\sqrt5$, so the correlation is $0.08 \times (25/8)/\sqrt5 = 0.112$
and $t = 0.112\sqrt{1\,255} = 3.96$, close to the 3.2 to 4.1 measured.
\end{solution}

\begin{solution}{exo:s1:cross-asset-lead-lag:3}
It sums products of returns over every pair of intervals that overlap, without sampling both series on a common grid; it neither discards ticks nor
creates the zero returns that previous-tick sampling does, which bias correlations toward zero (the Epps effect of Book~4).
\end{solution}

\begin{solution}{exo:s1:cross-asset-lead-lag:4}
$-0.51 + 1.0 = 0.49$ basis points before costs. The follower's observed price is its last tick, which arrives at random times: two seconds after the
leader's move, some of the follower's ticks have not yet printed the move, so a little of it is still ahead in the observed prices.
\end{solution}

\begin{solution}{exo:s1:cross-asset-lead-lag:5}
The window scan adds five lags' evidence into one statistic and pays for only 100 tests, so it has more power for a spread-out lead. The lag-by-lag
scan tests each lag alone against a stricter threshold, but a lag that happens to be strong in the sample (link~1 at five days, $t = 4.7$) can pass it
while its window statistic, diluted by weaker lags, does not.
\end{solution}

\begin{solution}{exo:s1:cross-asset-lead-lag:6}
It can report both: the correlated commodity inherits part of the true leader's predictive correlation. Regress the follower on both leaders' windows
together (the false one's coefficient should vanish), and prefer links with a mechanism.
\end{solution}

\begin{solution}{exo:s1:cross-asset-lead-lag:7}
The lag-by-lag scan finds two links (link~3 at one day and link~1 at five), the window scan all three planted links ($t$ of 5.3, 6.1 and 6.2) and a
fourth, commodity~10 to industry~4, which is false ($t = 3.7$): it rides partly on its correlation of 0.21 with the true leader of industry~4 and partly
on noise. The out-of-sample Sharpe ratio is 2.56 after costs; the false link costs a little.
\end{solution}

\begin{solution}{exo:s1:cross-asset-lead-lag:8}
30\,000 pair-lag tests with the 50 strongest selected: most are noise, and the backtest selected them in sample. Correct for the number of tests, require
a mechanism, estimate on one period and trade on another, and charge costs on 50 pairs traded daily.
\end{solution}

\begin{solution}{pb:s1:cross-asset-lead-lag:1}
\begin{enumerate}
\item A prediction of one market from another's past; information reaching prices at different speeds because investors specialise and attend to
few markets.
\item Industries such as retail, metal and petroleum forecast the market by up to two months; in the extended sample a robust core still does, but
fewer industries, with time variation.
\item Related supplier and customer industries cross-predict, less so when more informed investors follow them.
\item The big-to-small lead is mostly within industries and comes from slow adjustment to bad news.
\item 0.08 times the average of a commodity's last five daily returns, in three links; about 0.05 correlation per lag.
\item Lag by lag, 1\,000 tests at 4.06; the five-day window, 100 tests at 3.48.
\item One link at five days ($t = 4.7$); the other two links ($t = 4.1$ each), missing the first ($t = 3.2$); nothing false.
\item A Sharpe ratio of 2.09 before costs and 1.73 after.
\item A follower that sees the leader's price two seconds late, both with tick noise; Hayashi--Yoshida cross-correlation.
\item At two seconds, in all twenty sessions.
\item Trade the follower after a leader's move of more than 2 basis points in two seconds, hold two seconds, half a basis point each way.
\item 0.73, 0.35, $-0.51$ and $-0.81$ basis points at latencies of 0 to 3 seconds.
\item \textbf{Named result.} Daily: a five-day lead detected in two of three planted links by the window scan (one by the lag scan), traded at a Sharpe ratio
of 1.73 after costs out of sample; intraday: a two-second lead detected in all twenty sessions, worth 0.73 basis points per trade at no latency, 0.35 at
one second, and a loss from two seconds.
\item Statistics (few weak links among many candidates); speed and costs.
\item It sets the scan's shape and the number of tests, which decides the power.
\item Traders learn them and trade them away, as most published signals decay.
\item Rolling scans, split samples, and a link switched off in simulation to see how fast the test notices.
\item The futures-to-cash intraday lead.
\item Chapter~17 measured how much of a news move is left at each latency; here the lead of one asset over another is the news.
\item Whoever sees the leader first and trades the follower most cheaply.
\end{enumerate}
\end{solution}

\begin{solution}{iq:s1:cross-asset-lead-lag:1}
Correlate the follower's returns with the leader's past returns at the relevant horizons, in both directions, with an estimator suited to the data
(Hayashi--Yoshida for ticks); correct for the number of lags and pairs; check stability over time and out of sample; ask for a mechanism.
\end{solution}

\begin{solution}{iq:s1:cross-asset-lead-lag:2}
At 1\% about 100 of 10\,000 would be significant by chance; 40 is fewer than chance. Nothing has been found; apply a correction for multiple tests and a
mechanism before looking again.
\end{solution}

\begin{solution}{iq:s1:cross-asset-lead-lag:3}
Synchronise clocks (use exchange timestamps, or measure and correct each feed's delay), keep the ticks asynchronous and apply the Hayashi--Yoshida
cross-correlation over a grid of lags, and read the lead from the peak or the centre of symmetry of the correlation function.
\end{solution}

\begin{solution}{iq:s1:cross-asset-lead-lag:4}
Competitors got faster, the follower's liquidity providers started quoting off the leader, fees or spreads changed, or the firm's own latency grew;
measure the lead and the firm's latency again.
\end{solution}

\begin{solution}{iq:s1:cross-asset-lead-lag:5}
The link can break (a regime change, faster competitors), the follower can gap on its own news, and the position is unhedged against the follower's
specific risk; intraday, stale data or a slow feed turns the trade into adverse selection.
\end{solution}

\begin{solution}{iq:s1:cross-asset-lead-lag:6}
$\operatorname{Cov}(y_t, x_{t-k}) = \beta\sigma_x^2/L$, so $\rho_1 = \beta\sigma_x/(L\sigma_y)$. The window $W = \sum_k x_{t-k}$ has variance $L\sigma_x^2$ and
$\operatorname{Cov}(y_t, W) = \beta\sigma_x^2$, so $\rho_W = \beta\sigma_x/(\sqrt{L}\sigma_y)$. With the same sample, the ratio of t-statistics is
$\rho_W/\rho_1 = \sqrt{L}$: the window has $\sqrt{L}$ times the power of any single lag.
\end{solution}
